% Options for packages loaded elsewhere
\PassOptionsToPackage{unicode}{hyperref}
\PassOptionsToPackage{hyphens}{url}
\PassOptionsToPackage{dvipsnames,svgnames,x11names}{xcolor}
\documentclass[
  12pt,
  a4paper,
]{article}
\usepackage{xcolor}
\usepackage[margin=0.5in]{geometry}
\usepackage{amsmath,amssymb}
\setcounter{secnumdepth}{5}
\usepackage{iftex}
\ifPDFTeX
  \usepackage[T1]{fontenc}
  \usepackage[utf8]{inputenc}
  \usepackage{textcomp} % provide euro and other symbols
\else % if luatex or xetex
  \usepackage{unicode-math} % this also loads fontspec
  \defaultfontfeatures{Scale=MatchLowercase}
  \defaultfontfeatures[\rmfamily]{Ligatures=TeX,Scale=1}
\fi
\usepackage{lmodern}
\ifPDFTeX\else
  % xetex/luatex font selection
  \setmainfont[]{Times New Roman}
\fi
% Use upquote if available, for straight quotes in verbatim environments
\IfFileExists{upquote.sty}{\usepackage{upquote}}{}
\IfFileExists{microtype.sty}{% use microtype if available
  \usepackage[]{microtype}
  \UseMicrotypeSet[protrusion]{basicmath} % disable protrusion for tt fonts
}{}
\usepackage{setspace}
\makeatletter
\@ifundefined{KOMAClassName}{% if non-KOMA class
  \IfFileExists{parskip.sty}{%
    \usepackage{parskip}
  }{% else
    \setlength{\parindent}{0pt}
    \setlength{\parskip}{6pt plus 2pt minus 1pt}}
}{% if KOMA class
  \KOMAoptions{parskip=half}}
\makeatother
\usepackage{longtable,booktabs,array}
\newcounter{none} % for unnumbered tables
\usepackage{calc} % for calculating minipage widths
% Correct order of tables after \paragraph or \subparagraph
\usepackage{etoolbox}
\makeatletter
\patchcmd\longtable{\par}{\if@noskipsec\mbox{}\fi\par}{}{}
\makeatother
% Allow footnotes in longtable head/foot
\IfFileExists{footnotehyper.sty}{\usepackage{footnotehyper}}{\usepackage{footnote}}
\makesavenoteenv{longtable}
\usepackage{graphicx}
\makeatletter
\newsavebox\pandoc@box
\newcommand*\pandocbounded[1]{% scales image to fit in text height/width
  \sbox\pandoc@box{#1}%
  \Gscale@div\@tempa{\textheight}{\dimexpr\ht\pandoc@box+\dp\pandoc@box\relax}%
  \Gscale@div\@tempb{\linewidth}{\wd\pandoc@box}%
  \ifdim\@tempb\p@<\@tempa\p@\let\@tempa\@tempb\fi% select the smaller of both
  \ifdim\@tempa\p@<\p@\scalebox{\@tempa}{\usebox\pandoc@box}%
  \else\usebox{\pandoc@box}%
  \fi%
}
% Set default figure placement to htbp
\def\fps@figure{htbp}
\makeatother
% definitions for citeproc citations
\NewDocumentCommand\citeproctext{}{}
\NewDocumentCommand\citeproc{mm}{%
  \begingroup\def\citeproctext{#2}\cite{#1}\endgroup}
\makeatletter
 % allow citations to break across lines
 \let\@cite@ofmt\@firstofone
 % avoid brackets around text for \cite:
 \def\@biblabel#1{}
 \def\@cite#1#2{{#1\if@tempswa , #2\fi}}
\makeatother
\newlength{\cslhangindent}
\setlength{\cslhangindent}{1.5em}
\newlength{\csllabelwidth}
\setlength{\csllabelwidth}{3em}
\newenvironment{CSLReferences}[2] % #1 hanging-indent, #2 entry-spacing
 {\begin{list}{}{%
  \setlength{\itemindent}{0pt}
  \setlength{\leftmargin}{0pt}
  \setlength{\parsep}{0pt}
  % turn on hanging indent if param 1 is 1
  \ifodd #1
   \setlength{\leftmargin}{\cslhangindent}
   \setlength{\itemindent}{-1\cslhangindent}
  \fi
  % set entry spacing
  \setlength{\itemsep}{#2\baselineskip}}}
 {\end{list}}
\usepackage{calc}
\newcommand{\CSLBlock}[1]{\hfill\break\parbox[t]{\linewidth}{\strut\ignorespaces#1\strut}}
\newcommand{\CSLLeftMargin}[1]{\parbox[t]{\csllabelwidth}{\strut#1\strut}}
\newcommand{\CSLRightInline}[1]{\parbox[t]{\linewidth - \csllabelwidth}{\strut#1\strut}}
\newcommand{\CSLIndent}[1]{\hspace{\cslhangindent}#1}
\setlength{\emergencystretch}{3em} % prevent overfull lines
\providecommand{\tightlist}{%
  \setlength{\itemsep}{0pt}\setlength{\parskip}{0pt}}
\usepackage{pdflscape}
\usepackage{caption}
\captionsetup{margin=10pt,font=small,labelfont=bf}
%\usepackage{endfloat}

\usepackage{booktabs}
\usepackage{longtable}
\usepackage{array}
\usepackage{multirow}
\usepackage{wrapfig}
\usepackage{float}
\usepackage{colortbl}
\usepackage{pdflscape}
\usepackage{tabu}
\usepackage{threeparttable}
\usepackage{threeparttablex}
\usepackage[normalem]{ulem}
\usepackage{makecell}
\usepackage{xcolor}
\usepackage{bookmark}
\IfFileExists{xurl.sty}{\usepackage{xurl}}{} % add URL line breaks if available
\urlstyle{same}
\hypersetup{
  pdftitle={Fairness or Self-Interest? An Experiment on How Communication Shapes European Views of Fairness in the EU Asylum Relocation System},
  pdfauthor={(Alphabetical Order); Andrea Blasco; Begoña Cabeza Martínez; Rossella Icardi; Michal Krawczyk},
  colorlinks=true,
  linkcolor={Maroon},
  filecolor={Maroon},
  citecolor={blue},
  urlcolor={Blue},
  pdfcreator={LaTeX via pandoc}}

\title{Fairness or Self-Interest? An Experiment on How Communication Shapes European Views of Fairness in the EU Asylum Relocation System}
\author{(Alphabetical Order) \and Andrea Blasco \and Begoña Cabeza Martínez \and Rossella Icardi \and Michal Krawczyk}
\date{This version: Apr 30, 2026}

\begin{document}
\maketitle
\begin{abstract}
Despite intense political debate over reforming the EU asylum system, there is little evidence on what drives public support for alternative solutions. Using an experimental survey, we examine Europeans' fairness perceptions of three mechanisms for distributing asylum seekers among EU countries: no relocation, relocation proportional to GDP, and relocation proportional to the resident populations. By varying the information provided about the consequences of these mechanisms, we explored to what extent public perceptions are driven by concerns about outcomes versus principles of fairness. We find that most respondents (76\%) consider relocation to be the fairest option. However, there are large cross-country differences, with proportional relocation schemes perceived as fair by 58\% of respondents in Germany versus 87\% in Greece, which suggest a ``self-serving'' bias --- a tendency to rank as fairer the method that results in fewer asylum seekers for their own country. Moreover, when we communicate the expected number of asylum seekers for each relocation method, the self-serving bias intensifies significantly. By varying how the information is presented, we find a stronger effect when the communication displays \emph{absolute} figures instead of figures \emph{relative} to population. These findings shed new light on how public perceptions of fairness of the EU asylum system can be shaped by information.
\end{abstract}

\setstretch{1.1}
\clearpage

\section{Introduction}\label{introduction}

Over the past decade, the EU's asylum and migration system has been severely tested, from the 2015 escalation of the Syrian war to the recent conflict in Ukraine. These events exposed a disproportionate pressure placed on Member States at the EU's external borders. In response, the EU adopted the Pact on Migration and Asylum in May 2024, which establishes a permanent solidarity framework among Member States.

Under the Pact, the country of first entry remains responsible for asylum applications, while \emph{all} Member States commit to support the system through solidarity contributions, including the relocation of applicants or beneficiaries, financial contributions, or other operational support.\footnote{See European Commission (\citeproc{ref-eu2024pact}{2024})} However, there is no consensus among countries on how this responsibility should be shared. Because the implementation of solidarity mechanisms ultimately depends on political support at the national level, understanding public perceptions of fairness and solidarity in the EU asylum and migration system is crucial for ensuring the solidarity framework functions effectively, and maintaining social cohesion across the EU.

To explore these perceptions, we conducted an experimental survey in eight EU countries. Respondents were asked to evaluate the fairness of alternative asylum-seekers relocation systems under different information conditions about the consequences of each mechanism. Specifically, the survey asked respondents to rank the fairness of three mechanisms: no relocation, allocation proportional to each country's GDP, and allocation proportional to each country's population.

As in Bansak et al. (\citeproc{ref-bansak2017europeans}{2017}), we compared responses from those who received information about the consequences of implementing each method for the respondent's country, against those with no such information. Extending this approach, we additionally vary how consequences are communicated, presenting them either in \emph{absolute} terms or \emph{relative} to the country's population. This distinction is important because absolute figures may exaggerate perceived burdens, while relative measures better capture per-capita responsibility and may lead to different fairness judgments.

Our work builds on a large body of literature showing that people tend to combine several fairness principles, such as need, deservingness, or reciprocity, when judging distributional outcomes (\citeproc{ref-konow2003fairest}{Konow, 2003}). However, when one's own stakes are salient, a self-serving bias can shift those judgements so that the option that benefits agents the most appears ``fairest'' (\citeproc{ref-babcock1995biased}{Babcock et al., 1995}; \citeproc{ref-konow2000fair}{Konow, 2000}). This tension between ``distributional fairness'' and ``self-serving tendency'' appears particularly relevant when considering international solidarity in the asylum system. On the one hand, the principle of fairness suggests that countries should share migrants proportionally to their capacity. On the other hand, self-serving biases may lead to prioritising national interest. However, despite pioneering work by Bansak et al. (\citeproc{ref-bansak2017europeans}{2017}) to understand this issue, how these views shape public perceptions of fairness in the migration system is still poorly understood.

We find that nearly three-quarters of participants find some form of relocation most fair. However, there are important country differences. These are consistent with a self-serving bias: the tendency to judge the method that results in relatively few asylum seekers in one's own country as fairer. Moreover, communication strategies randomised as treatments significantly moderate this tendency. Explicitly spelling out the likely consequences of the implementation of each method in absolute terms (how many asylum seekers are expected in the participants' country per year) significantly strengthens the self-serving bias. The same information presented in relative terms (how many asylum seekers are to be expected per 100,000 inhabitants), showing that the number of asylum seekers is small compared to the entire population, has a weaker effect.

\section{Related Literature}\label{related-literature}

The EU Pact on Migration and Asylum is designed to manage the complex issue of migration more effectively. Its key component is the relocation system, aiming to create a fairer and more predictable process than the previous rules. The relocation of asylum seekers has been subject to increased controversy. This issue has fuelled significant political and social tensions for several reasons, including an uneven distribution whereby border countries bear most responsibility for processing asylum claims. Other EU member states were often unwilling to accept a share of these asylum seekers, leading to a major imbalance in burden-sharing (\citeproc{ref-thielemann2018refugee}{Thielemann, 2018}). The issue has become a powerful tool for populist and far-right political parties across Europe, using anti-immigrant rhetoric to mobilise voters and gain political power (\citeproc{ref-dinas2019waking}{Dinas et al., 2019}; \citeproc{ref-halla2017immigration}{Halla et al., 2017}).

Despite an intense political debate about asylum-seekers' relocation, there is still little evidence on Europeans' preferences. Survey evidence indicates that Europeans favour capacity-based relocation over both the Dublin ``first-entry'' rule and financial transfers. The seminal paper of Bansak et al. (\citeproc{ref-bansak2017europeans}{2017}), using a survey with 18,000 respondents across 15 EU+ countries, shows that 72\% of respondents backed proportional allocation even after learning it would raise their own country's intake. Similarly, a 26-country experiment by C. W. Cappelen et al. (\citeproc{ref-cappelen2023trade}{2023}) showed that most people prefer admitting asylum seekers to paying host states. A cross-national conjoint study in eight EU countries found that citizens favour refugee protection paired with clear limits or conditions, rather than pure first-entry allocations (\citeproc{ref-jeannet2021asylum}{Jeannet et al., 2021}). Existing research, hence, suggests that public preferences are more nuanced and conditional than commonly assumed. Recent UK-based findings by Corre and Tilley (\citeproc{ref-corre2024extent}{2024}) particularly highlight this complexity. They used a conjoint experiment to test whether current UK asylum policies align with public desire. Their results challenge the common belief that the public uniformly favours stricter policies than those currently in place. These results strongly suggest that to truly understand public preferences, researchers must disaggregate complex immigration policies into their specific components, moving beyond treating asylum as a single issue.

Citizens' widespread support for distributing asylum seekers based on a country's capacity implies a belief in the principle of proportional fairness and the equitable sharing of responsibility. However, this does not rule out that self-interest also shapes their attitudes, as some studies suggest. Bansak et al. (\citeproc{ref-bansak2017europeans}{2017}) found that the large majority of citizens favour an allocation policy proportional to each country's capacity over the status quo of allocation based on the country of first entry. This majority support, though weakened, persists even when respondents are informed that proportional allocation would increase the number of asylum seekers sent to their own country. These findings suggest that citizens care deeply about the fairness of the process.

Heizmann and Ziller (\citeproc{ref-heizmann2020willing}{2020}) analyses Eurobarometer data and finds that people living in countries that have taken in more asylum seekers and immigrants in the past are especially supportive of policies that redistribute wealth or resources, a preference they link to a mix of self-interest, concerns for their country's collective well-being, and moral values. Experimental work in Norway and Israel demonstrates that establishing clear proportional-allocation rules can mitigate NIMBYism (acronym for the ``not in my back yard'' attitude) by framing redistribution of refugees as a civic duty rather than an imposed burden (\citeproc{ref-solodoch2024overburdened}{Solodoch, 2024}). U.S. evidence from conjoint analysis further confirms that admission decisions are based more on collective good and ethical standards than on individual economic gain (\citeproc{ref-hainmueller2015hidden}{Hainmueller and Hopkins, 2015}).

Furthermore, a growing body of research confirms that public attitudes toward migration policy are malleable and can be meaningfully shifted through the use of concrete information and tailored framing. Studies show that specific informational framings can influence policy support: for example, informing U.S. respondents about the administrative burdens of legal immigration boosts support for more open policies (\citeproc{ref-kustov2025immigration}{Kustov and Landgrave, 2025}). Similarly, messaging congruent with personal values (using the Schwartz (\citeproc{ref-schwartz1992universals}{1992}) framework) proves more persuasive than generic humanitarian appeals in European campaigns (\citeproc{ref-dennison2020basic}{Dennison, 2020}), and correcting misperceptions about immigration's labour-market impact increases and sustains public support (\citeproc{ref-haaland2020labor}{Haaland and Roth, 2020}). In addition, narratives based on demographic challenges---such as focusing on labour shortages or an ageing population---can raise pro-immigration attitudes and political action (\citeproc{ref-facchini2022countering}{Facchini et al., 2022}). Conversely, Bansak et al. (\citeproc{ref-bansak2017europeans}{2017}) found that providing consequences-focused information (specifically, revealing that a proportional policy would increase asylum seeker numbers in a respondent's own country) appeared to boost self-interest, thereby slightly weakening support for the fairness-based proportional relocation in those countries where proportionality and self-interest don't align. A. W. Cappelen et al. (\citeproc{ref-cappelen2023second}{2023}) observed a self-serving bias where respondents deemed an allocation method fairer if it resulted in a lower national intake.

Drawing on the literature presented, we posit two core expectations (see the AEA RCT Registry for details). First, support for a proportional allocation system will be high: in each treatment, the share of participants considering the ``no relocation'' status quo as fairest will be significantly below the benchmark of one-third of all participants. Second, the perceived fairness of an option will negatively depend on the expected number of asylum seekers in one's own country under this option communicated to participants (self-serving bias). Specifically, consistent with research showing people overestimate migrant populations (\citeproc{ref-alesina2023immigration}{Alesina et al., 2023}), we expect self-serving bias to be more pronounced when the consequences are presented in absolute terms (thousands) compared to relative terms (hundreds per 100,000 inhabitants). This means that where a given relocation policy reduces the national burden with respect to the status quo (which is true for countries currently receiving a disproportionately high number of asylum applications), we predict that the policy will be rated as significantly fairer in the information treatments. Conversely, in countries in which a proportional policy increases the national burden, we predict that the policy will be rated as significantly less fair in the information treatments. In all cases, we expect that the detrimental or beneficial self-serving effect will be strongest when the information is presented in absolute terms rather than relative terms. This prediction is based on the literature suggesting that the same social effect presented in absolute rather than relative terms might be perceived as larger, affecting policy preferences (\citeproc{ref-pedersen2017ratio}{Pedersen, 2017}), a manifestation of the more general phenomenon of denominator neglect (\citeproc{ref-reyna2008numeracy}{Reyna and Brainerd, 2008}).

This work contributes to the literature on how public perceptions of fairness influence policy acceptance. Policy acceptance across all sectors depends on citizens viewing them as fair. Fairness judgments about who deserves what, often guided by deservingness heuristics, are a central component of distributive justice and shape how citizens evaluate who should contribute and who should be supported under a given policy (\citeproc{ref-almaas2010fairness}{Almås et al., 2010}; \citeproc{ref-cappelen2023trade}{C. W. Cappelen et al., 2023}; \citeproc{ref-corneo2002individual}{Corneo and Grüner, 2002}; \citeproc{ref-durante2014preferences}{Durante et al., 2014}; \citeproc{ref-fong2001social}{Fong, 2001}; \citeproc{ref-fong2011fairness}{Fong and Luttmer, 2011}; \citeproc{ref-van2006making}{Van Oorschot, 2006}, \citeproc{ref-van2000should}{2000}). Information can influence support, but its effects are heterogeneous: mobility information raises support (especially for opportunity policies), while inequality information mostly shifts concern and has large policy effects chiefly for the estate tax (\citeproc{ref-alesina2023immigration}{Alesina et al., 2023}; \citeproc{ref-kuziemko2015elastic}{Kuziemko et al., 2015}). In climate policy, support for carbon pricing rises when revenues are allocated in ways seen as fair, and acceptance can be improved with clear and trustworthy information on the policy's impact (\citeproc{ref-carattini2019win}{Carattini et al., 2019}; \citeproc{ref-maestre2019perceived}{Maestre-Andrés et al., 2019}).

Our study also adds to the literature in behavioural economics and decision making, which aims to understand how people respond to different ways of framing information. Work in medical research, for instance, has shown that people respond differently to risks when they are presented in absolute versus relative terms (\citeproc{ref-malenka1993framing}{Malenka et al., 1993}). Our work bridges these insights from medical and behavioural science literature and embeds them within a migration policy context.

\section{Methods and Data}\label{methods-and-data}

\subsection{Experimental design}\label{experimental-design}

A sample representative of the adult population (18 years and over) in eight EU countries (Bulgaria, France, Germany, Greece, Italy, Poland, Spain and Sweden) was drawn from Ipsos online panel by means of stratified random probability sampling. The stratification leveraged available profile data, such as respondents' age, gender and geographic region (NUTS 1), and pre-defined sub-sample sizes (quota) based on official population statistics provided by Eurostat. The overall sample size of the project included 8,122 participants. The main fieldwork took place between 30 September 2024 and 21 October 2024. Details of the field work and survey design are available elsewhere (\citeproc{ref-blasco2025fairness}{Blasco et al., 2025}).

The experiment tested how different types of information affected participants' ranking of the fairness of three relocation rules for asylum seekers across EU countries:\footnote{
  The experimental design and pre-analysis plan were pre-registered within the AEA RCT Registry (AEARCTR-0014194). The Ethics Research Board of the JRC approved the experimental protocol.}

\begin{itemize}
\item
  \textbf{No relocation} (the baseline): asylum seekers remain in the country where they first arrive;
\item
  Allocation \textbf{proportional to each country's population}: relocations ensure equity with population as a proxy for capacity;
\item
  Allocation \textbf{proportional to each country's GDP}: relocations ensure equity with GDP as a proxy for capacity.
\end{itemize}

Respondents were asked to rank these allocation rules from the fairest (rank = 1) to the least fair (rank = 3). Before they did so, we varied the format and amount of information available across three conditions:

\begin{itemize}
\tightlist
\item
  \textbf{No information} (the control),
\item
  \textbf{Absolute} information: the expected number of asylum applications under each relocation scheme,
\item
  \textbf{Relative} information: the expected number of asylum applications under each relocation scheme per 100,000 citizens.
\end{itemize}

Table \ref{tab:example} illustrates the experiment with a simplified scenario involving two countries, Country A and Country B.\footnote{
  This scenario is for illustrative purposes; it was not discussed with participants in the experiment.} Both countries have the same GDP, but Country A has more residents, while Country B received more than twice as many asylum applications, on average, in the past ten years (2014-2023). As a result, Country A will experience more asylum applications under proportional relocation than under no relocation. Furthermore, since countries are equally rich, but Country A is more densely populated, the difference in relocation outcomes will be greater under population-based relocation than under GDP-based relocation.

Under this scenario, study participants from Country A assigned to the absolute-information condition would be informed that their country would host 15,000 asylum seekers under no relocation, 30,000 under relocation by population, and 25,000 under relocation by GDP. In the relative information condition, the same information for Country A would be expressed per 100,000 citizens, with the provided values being 100, 200, and 167, respectively. This example illustrates that, while presenting information about asylum applications in absolute terms can involve large numbers, expressing them as relative figures often shows that the actual impact is modest. It also shows that for each relocation rule, some countries (here, both for relocation based on population and GDP, it is Country A) receive more asylum seekers than in the baseline. We will call such countries ``net receivers''. Remaining countries (here, Country B) receive fewer, becoming ``net senders''.

\begin{table}
\centering
\caption{\label{tab:example}Illustrative example for two countries}
\centering
\begin{tabu} to \linewidth {>{\raggedright\arraybackslash}p{10em}>{\raggedleft}X>{\raggedleft}X>{\raggedleft}X>{\raggedleft}X>{\raggedleft}X}
\toprule
\multicolumn{1}{c}{} & \multicolumn{2}{c}{Country A} & \multicolumn{2}{c}{Country B} & \multicolumn{1}{c}{All} \\
\cmidrule(l{3pt}r{3pt}){2-3} \cmidrule(l{3pt}r{3pt}){4-5} \cmidrule(l{3pt}r{3pt}){6-6}
 & Abs. & Rel. & Abs. & Rel. & \\
\midrule
\addlinespace[0.3em]
\multicolumn{6}{l}{\textbf{Hypothetical situation:}}\\
\hspace{1em}\em{Average yearly asylum applications (2014-23)} & 15.0K & 100 & 35.0K & 350 & 50.0K\\
\addlinespace\hline\hspace{1em}\em{Most recent population} & 15.0M & 100.0K & 10.0M & 100.0K & 25.0M\\
\addlinespace\hline\hspace{1em}\em{Most recent GDP} & 1.0B & 6.7M & 1.0B & 10.0M & 2.0B\\
\addlinespace\hline\addlinespace[0.3em]
\multicolumn{6}{l}{\textbf{Information treatment:}}\\
\hspace{1em}\em{No relocation} & 15.0K & 100 & 35.0K & 350 & 50.0K\\
\addlinespace\hline\hspace{1em}\em{Relocation by population} & 30.0K & 200 & 20.0K & 200 & 50.0K\\
\addlinespace\hline\hspace{1em}\em{Relocation by GDP} & 25.0K & 166.7 & 25.0K & 250 & 50.0K\\
\bottomrule
\end{tabu}
\end{table}

In Table \ref{tab:asylum-per-capita}, we report the actual figures for the expected asylum applications in each country under each information condition used in the experiment. They are based on population, GDP, and historical (2014-2023) asylum applications statistics.

\begin{table}
\centering
\caption{\label{tab:asylum-per-capita}Expected asylum seekers per 100 000 residents, under different allocation rules}
\centering
\begin{tabu} to \linewidth {>{\raggedright}X>{\raggedleft}X>{\raggedleft}X>{\raggedleft}X>{\raggedleft}X>{\raggedleft}X}
\toprule
Country & No relocation & Relocation by GDP & Relocation by pop. & Diff. (GDP) & Diff. (pop.)\\
\midrule
Greece & 397 & 95 & 169 & 302 & 228\\
Sweden & 344 & 230 & 169 & 114 & 175\\
Bulgaria & 178 & 65 & 169 & 113 & 9\\
Germany & 313 & 219 & 169 & 94 & 144\\
Spain & 138 & 136 & 169 & 2 & -31\\
Italy & 128 & 158 & 169 & -30 & -41\\
France & 151 & 184 & 169 & -33 & -18\\
Poland & 16 & 92 & 169 & -76 & -153\\
\bottomrule
\end{tabu}
\end{table}

The survey included several modules that asked respondents about their socio-demographics, perceptions of the EU asylum system, support for different migration policies, and other background questions. The full questionnaire is in the Appendix. A descriptive analysis of the responses is presented in the technical report by Blasco et al. (\citeproc{ref-blasco2025fairness}{2025}). In our analysis, we included control variables such as trust or political affiliation, as these are associated with attitudes towards migration (e.g., \citeproc{ref-canetti2016threatened}{Canetti et al., 2016}; \citeproc{ref-mitchell2021social}{Mitchell, 2021}). We also consider various socio-demographic factors, such as income and employment status, gender, age, and the parents' country of birth. Whether the respondent was born abroad may be relevant because foreigners may hold different opinions on migration issues than natives (\citeproc{ref-berlinschi2019migrants}{Berlinschi and Harutyunyan, 2019}).

\subsection{Data}\label{data}

Table \ref{tab:assignment-summary} in the Appendix illustrates the random assignment to the three information conditions in each country and overall. Each treatment group has about one third of total respondents per country, reflecting the stratified randomisation design.

Table \ref{tab:summary} illustrates the distribution of key pre-treatment characteristics, including sex, age, education, trust, and perceptions across the three treatment groups. On average, respondents were 50 years old (with 16 years SD), with about half (47\%) identifying as male. For simplicity of exposition and comparability across countries, education was merged into three broad categories: high (ISCED 5-8; 32\% of respondents), medium (ISCED 3-4; 48\%), and low (ISCED 0-2; 19\%). On measures of trust, respondents reported moderate levels of general trust in others (M = 4.5, SD = 2.7), trust in the EU (M = 4.3, SD = 2.7), and trust in government (M = 3.7, SD = 2.9).

Knowledge of the EU asylum system averaged 2.0 (SD = 0.7) on a 1-4 scale, thus the majority expresses no or basic knowledge. Most respondents reported there were ``too many'' asylum seekers in their country, with a mean of 2.6 on 1-3 scale from 1 = too few and 3 = too many. Similarly, on average, respondents view an unbalanced distribution of asylum seekers in the EU, reporting an average score of 3.1 (SD = 2.4) on 0-10 scale for the question about evenness across the EU. Yet, respondents reported broad consensus on the support on the importance of burden sharing among countries, providing equal rights, equal treatment, and having an equal distribution of asylum seekers relatively high, with means ranging from 6.9 to 7.4 on 0--10 scales.

Table \ref{tab:summary} also reports mean differences between the control group and the other treatment groups, with the standardised difference in parenthesis (the difference divided by the pooled standard deviation). For all these variables, the standardised differences between control and other groups is small, suggesting the distributions overlap well across treatment groups.\footnote{
  These differences were estimated using OLS controlling for country fixed effets, thus accounting for the stratified randomisation.}

\begin{table}
\centering
\caption{\label{tab:summary}Summary statistics (standardized mean difference in square brackets)}
\centering
\begin{tabu} to \linewidth {>{\raggedright\arraybackslash}p{15em}>{\raggedright}X>{\raggedright}X>{\raggedright}X}
\toprule
Variable & Control & Diff. Rel. & Diff. Abs.\\
\midrule
\addlinespace[0.3em]
\multicolumn{4}{l}{\textbf{Sociodemographics}}\\
\hspace{1em}Age (years) & 49.64 [3.01] & -0.28 [-0.02] & -0.15 [-0.01]\\
\hspace{1em}Male (\%) & 0.47 [0.94] & 0.01 [0.03] & 0.02 [0.03]\\
\hspace{1em}High education (\%) & 0.36 [0.76] & -0.01 [-0.02] & 0.01 [0.02]\\
\hspace{1em}Mid education (\%) & 0.48 [0.96] & 0.00 [0.00] & -0.01 [-0.01]\\
\hspace{1em}Low education (\%) & 0.15 [0.42] & 0.01 [0.03] & -0.00 [-0.00]\\
\addlinespace[0.3em]
\multicolumn{4}{l}{\textbf{Political attitudes}}\\
\hspace{1em}Trust others (0-10) & 4.47 [1.65] & 0.07 [0.03] & 0.09 [0.03]\\
\hspace{1em}Trust in government (0-10) & 3.67 [1.27] & 0.16 [0.06] & 0.17 [0.06]\\
\hspace{1em}Trust in the EU (0-10) & 4.30 [1.58] & 0.12 [0.04] & 0.12 [0.04]\\
\hspace{1em}Political right (0-10) & 5.28 [2.09] & -0.07 [-0.03] & -0.07 [-0.03]\\
\addlinespace[0.3em]
\multicolumn{4}{l}{\textbf{Other variables}}\\
\hspace{1em}Evenness of AS in EU (0-10) & 3.03 [1.25] & -0.10 [-0.04] & -0.01 [-0.01]\\
\hspace{1em}Value of burden sharing (0-10) & 7.44 [3.00] & -0.03 [-0.01] & 0.03 [0.01]\\
\hspace{1em}Value of equal distribution (0-10) & 6.92 [2.76] & -0.04 [-0.02] & -0.02 [-0.01]\\
\hspace{1em}Value of equal rights (0-10) & 6.88 [2.59] & -0.01 [-0.00] & 0.07 [0.03]\\
\hspace{1em}Value of equal treatment (0-10) & 7.14 [2.74] & -0.07 [-0.03] & 0.04 [0.01]\\
\hspace{1em}Value of quick process (0-10) & 6.96 [2.66] & -0.09 [-0.04] & 0.02 [0.01]\\
\hspace{1em}Objective Asylum decisions (0-10) & 5.30 [1.95] & -0.06 [-0.02] & -0.11 [-0.04]\\
\hspace{1em}Frauds detection ability (0-10) & 3.86 [1.40] & 0.09 [0.03] & 0.12 [0.04]\\
\hspace{1em}Share legal migrants & 45.49 [1.81] & -0.20 [-0.01] & 0.81 [0.03]\\
\hspace{1em}Time to complete (mins) & 21.72 [1.11] & -1.11 [-0.06] & -0.25 [-0.01]\\
\bottomrule
\end{tabu}
\end{table}

\subsection{Fair share and political leanings}\label{fair-share-and-political-leanings}

\begin{itemize}
\tightlist
\item
  TODO: I can look at measures of fairness self-reported correlated with left-right
\item
  TODO: Left-right vs fairness. Fair share of migrants.
\end{itemize}

\section{Results}\label{results}

\subsection{Perceived fairness of the relocation systems}\label{perceived-fairness-of-the-relocation-systems}

Figure \ref{fig:toprank} presents the share of respondents ranking the fairest asylum-seeker allocation mechanism across countries and information conditions. Overall, relocation mechanisms are ranked as fairer than no relocation (70\% or higher share of respondents), with allocation based on GDP receiving the highest share in the no information group (39\%), followed closely by allocation based on population (35\%). In the control group, only 26\% considered no relocation at all to be the fairest option. A one-sample proportion test shows that this result is significantly lower (p\textless0.01) than the 33\% expected if individuals were ranking alternatives at random or were indifferent among the options.

Figure \ref{fig:toprank} also reveals considerable variation across countries in the control group. In Germany, only 58\% ranked some form of relocation as the fairest option, which is the lowest share among the selected countries, while in Greece the proportion was highest at 87\%. Relocation by GDP was commonly perceived as the fairest in Bulgaria (53\%), Greece (48\%) and Spain (48\%), but not in Germany (19\%) or Sweden (27\%). Views on relocation by population also varied: the share of respondents assigning it the top rank range from 24\% in Poland to 45\% in Italy. These stark differences suggest that country-specific conditions, such as the political climate, refugee situation, and socio-economic differences, play a significant role in shaping public perceptions of fairness.

As illustrated by Figure \ref{fig:toprank}, providing information (absolute or relative) tends to increase support for relocation compared with having no information, highlighting the role of information in shaping public preferences, as we discuss next.

\begin{figure}
\includegraphics[width=\textwidth]{/Users/mrb/Research/migrant_relocation_fairness/draft/jebo-rev1/main_files/figure-latex/toprank-1} \caption[Fairest asylum-seeker allocation mechanism.]{Fairest asylum-seeker allocation mechanism. Share of respondents selecting the top-ranked relocation mechanism across information conditions}\label{fig:toprank}
\end{figure}

\subsection{Average Treatment Effects on Fairness Rankings}\label{average-treatment-effects-on-fairness-rankings}

Figure \ref{fig:mean-rank} shows mean differences in fairness rankings across countries for the three allocation schemes, comparing respondents who received absolute or relative information about relocated asylum seekers with those who received no information. In countries that would host fewer asylum seekers under proportional relocation (``net senders''), such as Bulgaria, Germany, Greece, and Sweden, the no-relocation option receives higher fairness rankings in the absolute-information condition, indicating stronger perceived fairness. This effect is largest in Germany and Sweden (exceeding 0.4 points) and is statistically significant (p \textless{} 0.05) in all countries except Bulgaria. By contrast, in countries that would receive fewer asylum seekers under no relocation (``net receivers''), including France, Italy, Poland, and Spain, providing this information lowers the average rank of the no-relocation option by about 0.3 points, indicating weaker perceived fairness.

Figure \ref{fig:mean-rank} also reveals remarkable differences in fairness rankings between the absolute and relative information conditions. In Spain, Poland, and Italy, for example, the difference fairness rankings between absolute information and no information is significantly larger than under relative information. Overall, these results confirm our hypotheses that individuals react more when information is presented in absolute than in relative terms.\footnote{As a formal test of the overall associations between fairness rankings and treatment assignment, we used separate Chi-squared tests within each country with Bonferroni correction for multiple testing. Even after applying Bonferroni correction, all the country-specific associations were statistically significant (p.~\textless{} 0.01), confirming a robust treatment effect in most countries, as hypothesised.}

\begin{figure}
\includegraphics[width=\textwidth]{/Users/mrb/Research/migrant_relocation_fairness/draft/jebo-rev1/main_files/figure-latex/mean-rank-1} \caption[Mean differences in fairness rankings]{Information provision and fairness rankings. Points show mean differences in fairness rankings (with 95\% confidence intervals) between individuals receiving absolute or relative information and those receiving no information. Lower ranks indicate higher perceived fairness. Countries are classified as net receivers if expected asylum applications are higher under relocation than under no relocation (based on GDP or population), and as net senders otherwise.}\label{fig:mean-rank}
\end{figure}

\subsection{The role of asylum seekers}\label{the-role-of-asylum-seekers}

We further investigate the role of the expected asylum seekers under each relocation scheme using a rank-ordered logit model for the probability of ranking each relocation scheme higher in the fairness ranking (see the Appendix).

The baseline specification of the rank-ordered logit model, which assumes that the effect of asylum seekers does not vary across treatments, indicates that proportional schemes were generally ranked as fairer, holding constant the number of asylum seekers (see Figure \ref{fig:rologit-base} in the Appendix). This is evidenced by the positive and statistically significant intercept interactions for GDP (0.64, p \textless{} 0.05) and population (0.57, p \textless{} 0.05). In contrast, the negative coefficient for asylum seekers (−0.35, p \textless{} 0.05) suggests that schemes allocating higher numbers of asylum seekers were perceived as less fair. These results are consistent with our hypotheses and earlier analyses.\footnote{
  In this specification, we do not expect strong average treatment effects, as treatment effects should have opposite signs across countries, offsetting one another. Among the interaction terms, however, the effect of relocation by population under the control condition (0.13, p \textless{} 0.05) is positive and statistically significant. This indicates that fairness perceptions were, on average, responsive to information about differences in asylum seeker allocations, whereas the remaining interaction effects are not statistically significant.}

\begin{table}
\centering
\caption{\label{tab:coeffs}Regression coefficients from rank-ordered logit on fairness rankings}
\centering
\begin{tabu} to \linewidth {>{\raggedright\arraybackslash}p{15em}>{\centering}X>{\centering}X>{\centering}X}
\toprule
Coefficient & No info & Absolute & Relative\\
\midrule
Relocation by GDP & 0.74 & 0.56 & 0.67\\
 & (0.04) & (0.04) & \vphantom{1} (0.04)\\
 & {}[0.7,0.8] & {}[0.5,0.6] & \vphantom{1} {}[0.6,0.7]\\
Relocation by population & 0.75 & 0.54 & 0.66\\
 & (0.04) & (0.04) & (0.04)\\
\addlinespace
 & {}[0.7,0.8] & {}[0.5,0.6] & {}[0.6,0.7]\\
Asylum seekers (std.) & -0.13 & -0.57 & -0.40\\
 & (0.02) & (0.03) & (0.03)\\
 & {}[-0.2,-0.1] & {}[-0.6,-0.5] & {}[-0.4,-0.3]\\
\bottomrule
\multicolumn{4}{l}{\rule{0pt}{1em}\textit{Note: }}\\
\multicolumn{4}{l}{\rule{0pt}{1em}Standard errors in parenthesis, and 95\% confidence intervals in square brackets.}\\
\end{tabu}
\end{table}

Figure \ref{fig:rologit-ate} and Table \ref{tab:coeffs} report the results of the rank-ordered logit model estimated separately for each information condition, thus all coefficients to vary across the treatment conditions. Across all groups, alternatives with proportional relocation, either by GDP or by populations, were consistently ranked as fairer, with positive and significant coefficients in each condition, again confirming earlier results.

However, the strength of these effects varied: the control group showed the largest by-GDP and by-population effects (0.74 and 0.75, respectively), while these effects were smaller under the absolute and relative information conditions. The negative coefficients for asylum seekers indicate that schemes allocating more asylum seekers were viewed as less fair, though the magnitude of this effect differed significantly across groups. It was strongest in the absolute-information condition (−0.57, 95\% CI: -0.62, -0.52), moderate in the relative-information condition (−0.40, 95\% CI: -.45, -.35), and weakest in the control group (−0.13, 95\% CI: -.17, -0.08). This pattern confirms that providing information, particularly absolute information, heightened sensitivity to the number of asylum seekers when evaluating fairness.

\begin{figure}
\includegraphics[width=\textwidth]{/Users/mrb/Research/migrant_relocation_fairness/draft/jebo-rev1/main_files/figure-latex/rologit-ate-1} \caption{Estimates and 95\% CIs from separate rank-ordered logit regressions on the probability of ranking an alternative relocation scheme as fairer, under each information treatment. $\beta_{AS}^{\tau}$ represents the effect of a SD change in expected asylum seekers on the probability of ranking an allocation scheme as fairer, under treatment $\tau$. $\alpha_{j}^{\tau}$ represents the effect of a proportional allocation mechanism $j$ relative to the no relocation alternative (omitted), under treatment $\tau$. This shows that the effects associated with a change in expected asylum seekers (AS) are stronger under absolute than relative information. A change in AS has a significant effect on fairness rankings even under no information, indicating some awareness of respondents beyond our intervention.}\label{fig:rologit-ate}
\end{figure}

\subsection{Heterogenous effects}\label{heterogenous-effects}

\subsubsection{Political orientation}\label{political-orientation}

Figure \ref{fig:political} shows that the effect of information on perceived fairness is strongly moderated by political orientation. The figure reports intercepts and slopes from the rank-ordered logit model, stratified by respondents' political orientation (discretised into left, neutral, and right).\footnote{
  Since we centered the expected number of asylum seekers for each alternative, the model's intercepts for the type of relocation (by GDP or by population) represent the average treatment effect on the probability of ranking the scheme as fairer than having no relocation at all.} Across all groups, respondents assign a positive ``premium'' to relocation schemes relative to no relocation: intercepts are significantly above zero in all conditions. This premium is largest among left-leaning respondents, smaller among centrists, and weakest among right-leaning respondents. Intercepts are largely similar across information treatments.

By contrast, information substantially shapes sensitivity to the number of asylum seekers. In all models, alternatives with higher numbers of asylum seekers are associated with lower perceived fairness, but both absolute and relative information amplify this effect, especially among centrist and right-leaning respondents. The partisan divide is most pronounced under absolute information: right-leaning respondents react most negatively to higher numbers (−0.76), whereas left-leaning respondents are considerably less sensitive (−0.35), indicating a sharp ideological gap in fairness perceptions.

\begin{figure}
\includegraphics[width=\textwidth]{/Users/mrb/Research/migrant_relocation_fairness/draft/jebo-rev1/main_files/figure-latex/political-1} \caption{Heterogeneous Treatment Effects by Political orientation. Coefficients and 95\% confidence intervals from rank-ordered regression on the probability of ranking an alternative relocation shceme as fairer, stratified by the self-reported political orientation.}\label{fig:political}
\end{figure}

\subsubsection{Fair share}\label{fair-share}

Figure \ref{fig:fair} illustrates the heteogeneity of the effects between individuals who reported the number of migrants in their own country as ``fair'' or ``too low'' vs those perceiving it as ``too high''. Both groups show higher perceived fairness for relocation schemes compared to no relocation, with all intercepts significantly positive. However, among respondents reporting that the current share is ``fair or too low,'' we do not detect a meaningful average treatment effect of framing: support is lowest under the Absolute condition, particularly for relocation by population (Estimate = 0.56, SE = 0.08) compared with GDP-based relocation (Estimate = 0.73, SE = 0.08). By contrast, in the group perceiving the share as ``too high,'' the Absolute framing significantly reduced support relative to Control for both GDP and population schemes (GDP: 0.29 vs.~0.61; Population: 0.38 vs.~0.65), with Relative framing falling in between and not significantly different from Control.

At the same time, for respondents reporting a ``too high'' share, we observe strong and significant heterogeneous effects of the expected number of asylum seekers on each relocation scheme. The negative impact of a one standard deviation increase in asylum seekers was largest under the Absolute condition (Estimate = -0.79, SE = 0.04), followed by the Relative condition (-0.56, SE = 0.04), and smallest in the Control group (-0.19, SE = 0.04). This corresponds to a difference of approximately 0.60 between Absolute and Control, substantially larger than the corresponding difference in the ``fair or too low'' group, which is 0.26.

\begin{figure}
\centering
\pandocbounded{\includegraphics[keepaspectratio,alt={\label{fig:fair}Heterogeneous Treatment Effects by Perceived Fairness of Migrants' Distribution in the EU. Coefficients and 95\% confidence intervals from rank-ordered regression on the probability of ranking an alternative relocation shceme as fairer, stratified by whether respondents perceive the share of refugees in their country as `too high' or not.}]{/Users/mrb/Research/migrant_relocation_fairness/draft/jebo-rev1/main_files/figure-latex/fair-1.pdf}}
\caption{\label{fig:fair}Heterogeneous Treatment Effects by Perceived Fairness of Migrants' Distribution in the EU. Coefficients and 95\% confidence intervals from rank-ordered regression on the probability of ranking an alternative relocation shceme as fairer, stratified by whether respondents perceive the share of refugees in their country as `too high' or not.}
\end{figure}

\subsubsection{Age}\label{age}

Figure \ref{fig:age} illustrates the heterogeneity of the effects between younger and older respondents. Both age groups show higher support for relocation schemes compared to no relocation, with all intercepts significantly positive. However, among younger respondents, we do not detect a meaningful average treatment effect of information provision, which is instead notable for older respondents (45 years old and above).

\begin{figure}
\centering
\pandocbounded{\includegraphics[keepaspectratio,alt={\label{fig:age}Heterogeneous Treatment Effects by Age. Coefficients and 95\% confidence intervals from rank-ordered regression on the probability of ranking an alternative relocation shceme as fairer, stratified by the respondents' age.}]{/Users/mrb/Research/migrant_relocation_fairness/draft/jebo-rev1/main_files/figure-latex/age-1.pdf}}
\caption{\label{fig:age}Heterogeneous Treatment Effects by Age. Coefficients and 95\% confidence intervals from rank-ordered regression on the probability of ranking an alternative relocation shceme as fairer, stratified by the respondents' age.}
\end{figure}

\subsubsection{Konowledge of the Asylum System}\label{konowledge-of-the-asylum-system}

Figure \ref{fig:age} illustrates the heterogeneity of effects stratified by self-reported knowledge or awareness of the rules governing the EU asylum system. Remarkably, we do not observe significant differences between individuals with basic knowledge and those with more advanced understanding.

\begin{figure}
\centering
\pandocbounded{\includegraphics[keepaspectratio,alt={\label{fig:knowledge}Heterogeneous Treatment Effects by Knowledge of the Asylum System. Coefficients and 95\% confidence intervals from rank-ordered regression on the probability of ranking an alternative relocation shceme as fairer, stratified by the self-reported knowledge or awareness of the rules governing the asylum system in the EU.}]{/Users/mrb/Research/migrant_relocation_fairness/draft/jebo-rev1/main_files/figure-latex/knowledge-1.pdf}}
\caption{\label{fig:knowledge}Heterogeneous Treatment Effects by Knowledge of the Asylum System. Coefficients and 95\% confidence intervals from rank-ordered regression on the probability of ranking an alternative relocation shceme as fairer, stratified by the self-reported knowledge or awareness of the rules governing the asylum system in the EU.}
\end{figure}

\subsubsection{Trust in the EU}\label{trust-in-the-eu}

Figure \ref{fig:trust} illustrates the heterogeneity of effects stratified by individuals' overall trust in the EU. We find that sensitivity to information is particularly pronounced among those with low levels of trust.

\begin{figure}
\includegraphics[width=\textwidth]{/Users/mrb/Research/migrant_relocation_fairness/draft/jebo-rev1/main_files/figure-latex/trust-1} \caption{Heterogeneous Treatment Effects by Trust in the EU. Coefficients and 95\% confidence intervals from rank-ordered regression on the probability of ranking an alternative relocation shceme as fairer, stratified by the overall trust in the EU. The original 0-10 scale was grouped into high, medium, and low categories.}\label{fig:trust}
\end{figure}

\section{Discussion}\label{discussion}

This study used an experimental survey to examine Europeans' perceptions of the fairness of different mechanisms for distributing asylum applications among EU countries. By varying the information provided about the consequences of these mechanisms, it explored the extent to which public perceptions of fairness are driven by concerns about outcomes versus principles of fairness.

Consistent with prior research (\citeproc{ref-bansak2017europeans}{Bansak et al., 2017}), we find that a substantial majority (76\%) of respondents across eight EU countries view some form of proportional relocation of applicants as \emph{fairer} than no relocation at all. Moreover, most people view proportional relocation systems as \emph{fairer} even when it would result in their own country receiving more asylum applications, as the share of those ranking relocation as fairer than no relocation exceeds 50\% in all countries.

Yet, we find substantial variation in fairness perceptions across countries. In the control group, where no information about the outcomes was provided, we find that between 58\% respondents in Germany to 87\% respondents in Greece view some form of within-EU relocation as fairer than no relocation at all. This gap could be partly explained by concerns about the outcomes. For instance, a border country like Greece would experience more applications under a no relocation system, while Germany would experience fewer applications. This finding underscores the importance of disentangling concerns about outcomes from fairness.

Bansak et al. (\citeproc{ref-bansak2017europeans}{2017}) finds that varying the information available about outcomes significantly affects public support for different relocation systems. Building on this insight, our research demonstrates that providing information about the consequences is key in shaping public perceptions of the fairness of these systems. Thus, our study establishes a direct link between political support identified by Bansak et al. (\citeproc{ref-bansak2017europeans}{2017}) and the perceived fairness of different policy options. As such, our study adds to a growing literature exploring the role of public perceptions of fairness in regulation, such as Brick and Visser (\citeproc{ref-brick2015fair}{2015}) for climate regulation.

By providing information about the consequences of the allocation schemes, we find that when individuals learn about the outcomes, the initial preference for proportional schemes is offset if the relocation mechanism results in one standard deviation more asylum applications. In other words, Europeans are more likely to perceive proportional allocation systems as fair when the increase in the number of asylum seekers implied by these systems is not too large. A first key takeaway is that information matters: receiving information about the consequences of proportional mechanisms may alter perceptions of fairness, ultimately shaping public support for relocation policies.

Another key implication is that framing matters. Differences that appear large in absolute terms may seem small when presented in relative terms. Indeed, in our randomised experiment, we find that expressing the outcomes of different relocation systems in absolute numbers leads some individuals to react more strongly to the information. When no information is provided, a regime that increases expected asylum applications by one standard deviation is 13\% less likely to be seen as fair. Once information is revealed, fairness drops by 40\% with relative information and by 56\% with absolute information. Thus, providing absolute information produces an additional 16-percentage point decrease in perceived fairness compared to relative information. This finding highlights another way in which public opinion can be influenced by the framing of identical information.

Our analysis of heterogeneous effects provides additional evidence on the mechanisms through which absolute vs relative information works. These differing impacts are not uniform across the population but are highly dependent on individual factors, including political orientation and demographics. Across most education groups, we find a consistently negative and statistically significant correlation between the number of asylum seekers and perceived fairness. The magnitude of this association was larger in the absolute-information group, moderate in the relative-information group, and modest in the control group. However, the effect of asylum seekers was weakest for low-educated respondents in all groups. These findings suggest a minor sensitivity of low-educated individuals to information. With regard to political orientation, both GDP- and population-based relocation schemes were consistently viewed as fairer than no relocation, with the largest effects observed for left-leaning respondents under absolute and control conditions.

Our research relates to conspicuous experimental literature in other fields that has shown the sensitivity of human decision-making to absolute vs relative values (\citeproc{ref-kratochvil2024fallacy}{Kratochvı́l and Havlı́ček, 2024}; \citeproc{ref-malenka1993framing}{Malenka et al., 1993}). In medical research, for example, effect sizes or risk reductions of interventions may appear very large or quite small, depending on the use of relative or absolute measures on the same dataset (\citeproc{ref-malenka1993framing}{Malenka et al., 1993}). The use of relative risk reduction has been shown to be more likely to influence the decision to accept a treatment than absolute risk reduction. In political communication, the main reason why relative values are so widespread is that absolute values are mostly not of much use for comparison among countries because they vary in population size. It is worth noticing, however, that the use of relative information may lead to misleading conclusions (\citeproc{ref-silva2020per}{Silva, 2020}). They should, therefore, be used when the measure of interest increases proportionally with population size. Our study provides additional evidence on the effect of relative vs absolute communication in the field of migration policies.

This study has some key limitations. As the evidence is survey-based, the findings rely on self-reported attitudes and are susceptible to biases such as social desirability or a gap between stated and actual behaviour. The scope of the research is also geographically constrained: drawing data from eight countries limits the generalizability of the results to global or even broader regional contexts. Finally, the use of hypothetical scenarios, though necessary for experimental control, means that the observed reactions to absolute versus relative information may not perfectly translate to real-world policy decisions.

\section*{References}\label{references}
\addcontentsline{toc}{section}{References}

\protect\phantomsection\label{refs}
\begin{CSLReferences}{1}{0}
\bibitem[\citeproctext]{ref-alesina2023immigration}
Alesina, A., Miano, A., Stantcheva, S., 2023. Immigration and redistribution. The Review of Economic Studies 90, 1--39.

\bibitem[\citeproctext]{ref-almaas2010fairness}
Almås, I., Cappelen, A.W., Sørensen, E.Ø., Tungodden, B., 2010. Fairness and the development of inequality acceptance. Science 328, 1176--1178.

\bibitem[\citeproctext]{ref-babcock1995biased}
Babcock, L., Loewenstein, G., Issacharoff, S., Camerer, C., 1995. Biased judgments of fairness in bargaining. The American Economic Review 85, 1337--1343.

\bibitem[\citeproctext]{ref-bansak2017europeans}
Bansak, K., Hainmueller, J., Hangartner, D., 2017. Europeans support a proportional allocation of asylum seekers. Nature Human Behaviour 1, 0133.

\bibitem[\citeproctext]{ref-berlinschi2019migrants}
Berlinschi, R., Harutyunyan, A., 2019. Do migrants think differently? Evidence from eastern european and post-soviet states. International Migration Review 53, 831--868.

\bibitem[\citeproctext]{ref-blasco2025fairness}
Blasco, A., Cabeza Martinez, B., Icardi, R., Krawczyk, M., Seiger, F., 2025. Public perceptions of fairness in the european migration and asylum system. Publications Office of the European Union, Luxembourg. \url{https://doi.org/10.2760/1271462}

\bibitem[\citeproctext]{ref-brick2015fair}
Brick, K., Visser, M., 2015. What is fair? An experimental guide to climate negotiations. European Economic Review 74, 79--95.

\bibitem[\citeproctext]{ref-canetti2016threatened}
Canetti, D., Snider, K.L., Pedersen, A., Hall, B.J., 2016. Threatened or threatening? How ideology shapes asylum seekers' immigration policy attitudes in israel and australia. Journal of refugee studies 29, 583--606.

\bibitem[\citeproctext]{ref-cappelen2023second}
Cappelen, A.W., Cappelen, C., Tungodden, B., 2023. Second-best fairness: The trade-off between false positives and false negatives. American Economic Review 113, 2458--2485.

\bibitem[\citeproctext]{ref-cappelen2023trade}
Cappelen, C.W., Sicakkan, H.G., Van Wolleghem, P.G., 2023. The trade-off between admitting and paying: Experimental evidence on attitudes towards asylum responsibility-sharing. European Union Politics 24, 470--493.

\bibitem[\citeproctext]{ref-carattini2019win}
Carattini, S., Kallbekken, S., Orlov, A., 2019. How to win public support for a global carbon tax. Nature 565, 289--291.

\bibitem[\citeproctext]{ref-corneo2002individual}
Corneo, G., Grüner, H.P., 2002. Individual preferences for political redistribution. Journal of public Economics 83, 83--107.

\bibitem[\citeproctext]{ref-corre2024extent}
Corre, T.L., Tilley, J., 2024. To what extent does asylum policy match public policy preferences? International Migration Review 01979183241253502.

\bibitem[\citeproctext]{ref-dennison2020basic}
Dennison, J., 2020. A basic human values approach to migration policy communication. Data \& Policy 2, e18.

\bibitem[\citeproctext]{ref-dinas2019waking}
Dinas, E., Matakos, K., Xefteris, D., Hangartner, D., 2019. Waking up the golden dawn: Does exposure to the refugee crisis increase support for extreme-right parties? Political analysis 27, 244--254.

\bibitem[\citeproctext]{ref-durante2014preferences}
Durante, R., Putterman, L., Van der Weele, J., 2014. Preferences for redistribution and perception of fairness: An experimental study. Journal of the European Economic Association 12, 1059--1086.

\bibitem[\citeproctext]{ref-eu2024pact}
European Commission, 2024. \href{https://op.europa.eu/en/publication-detail/-/publication/a426d987-0767-11ef-a251-01aa75ed71a1/}{{Pact on migration and asylum}}.

\bibitem[\citeproctext]{ref-facchini2022countering}
Facchini, G., Margalit, Y., Nakata, H., 2022. Countering public opposition to immigration: The impact of information campaigns. European Economic Review 141, 103959.

\bibitem[\citeproctext]{ref-fong2001social}
Fong, C., 2001. Social preferences, self-interest, and the demand for redistribution. Journal of Public economics 82, 225--246.

\bibitem[\citeproctext]{ref-fong2011fairness}
Fong, C.M., Luttmer, E.F., 2011. Do fairness and race matter in generosity? Evidence from a nationally representative charity experiment. Journal of Public Economics 95, 372--394.

\bibitem[\citeproctext]{ref-haaland2020labor}
Haaland, I., Roth, C., 2020. Labor market concerns and support for immigration. Journal of Public Economics 191, 104256.

\bibitem[\citeproctext]{ref-hainmueller2015hidden}
Hainmueller, J., Hopkins, D.J., 2015. The hidden american immigration consensus: A conjoint analysis of attitudes toward immigrants. American journal of political science 59, 529--548.

\bibitem[\citeproctext]{ref-halla2017immigration}
Halla, M., Wagner, A.F., Zweimüller, J., 2017. Immigration and voting for the far right. Journal of the European economic association 15, 1341--1385.

\bibitem[\citeproctext]{ref-hausman1987specifying}
Hausman, J.A., Ruud, P.A., 1987. Specifying and testing econometric models for rank-ordered data. Journal of econometrics 34, 83--104.

\bibitem[\citeproctext]{ref-heizmann2020willing}
Heizmann, B., Ziller, C., 2020. Who is willing to share the burden? Attitudes towards the allocation of asylum seekers in comparative perspective. Social Forces 98, 1026--1051.

\bibitem[\citeproctext]{ref-jeannet2021asylum}
Jeannet, A.-M., Heidland, T., Ruhs, M., 2021. What asylum and refugee policies do europeans want? Evidence from a cross-national conjoint experiment. European Union Politics 22, 353--376.

\bibitem[\citeproctext]{ref-konow2003fairest}
Konow, J., 2003. Which is the fairest one of all? A positive analysis of justice theories. Journal of economic literature 41, 1188--1239.

\bibitem[\citeproctext]{ref-konow2000fair}
Konow, J., 2000. Fair shares: Accountability and cognitive dissonance in allocation decisions. American economic review 90, 1072--1092.

\bibitem[\citeproctext]{ref-kratochvil2024fallacy}
Kratochvı́l, L., Havlı́ček, J., 2024. The fallacy of global comparisons based on per capita measures. Royal Society Open Science 11, 230832.

\bibitem[\citeproctext]{ref-kustov2025immigration}
Kustov, A., Landgrave, M., 2025. Immigration is difficult?! Informing voters about immigration policy fosters pro-immigration views. Journal of Experimental Political Science 1--13.

\bibitem[\citeproctext]{ref-kuziemko2015elastic}
Kuziemko, I., Norton, M.I., Saez, E., Stantcheva, S., 2015. How elastic are preferences for redistribution? Evidence from randomized survey experiments. American Economic Review 105, 1478--1508.

\bibitem[\citeproctext]{ref-maestre2019perceived}
Maestre-Andrés, S., Drews, S., Van Den Bergh, J., 2019. Perceived fairness and public acceptability of carbon pricing: A review of the literature. Climate policy 19, 1186--1204.

\bibitem[\citeproctext]{ref-malenka1993framing}
Malenka, D.J., Baron, J.A., Johansen, S., Wahrenberger, J.W., Ross, J.M., 1993. The framing effect of relative and absolute risk. Journal of general internal medicine 8, 543--548.

\bibitem[\citeproctext]{ref-mitchell2021social}
Mitchell, J., 2021. Social trust and anti-immigrant attitudes in europe: A longitudinal multi-level analysis. Frontiers in sociology 6, 604884.

\bibitem[\citeproctext]{ref-pedersen2017ratio}
Pedersen, R.T., 2017. Ratio bias and policy preferences: How equivalency framing of numbers can affect attitudes. Political Psychology 38, 1103--1120.

\bibitem[\citeproctext]{ref-reyna2008numeracy}
Reyna, V.F., Brainerd, C.J., 2008. Numeracy, ratio bias, and denominator neglect in judgments of risk and probability. Learning and individual differences 18, 89--107.

\bibitem[\citeproctext]{ref-schwartz1992universals}
Schwartz, S.H., 1992. Universals in the content and structure of values: Theoretical advances and empirical tests in 20 countries, in: Advances in Experimental Social Psychology. Elsevier, pp. 1--65.

\bibitem[\citeproctext]{ref-silva2020per}
Silva, W., 2020. Per capita death and infection rates should be avoided in international comparisons. Public Health 186, 18.

\bibitem[\citeproctext]{ref-solodoch2024overburdened}
Solodoch, O., 2024. Overburdened? How refugee dispersal policies can mitigate NIMBYism and public backlash. Journal of European Public Policy 31, 2455--2482.

\bibitem[\citeproctext]{ref-thielemann2018refugee}
Thielemann, E., 2018. Why refugee burden-sharing initiatives fail: Public goods, free-riding and symbolic solidarity in the EU. JCMS: Journal of Common Market Studies 56, 63--82.

\bibitem[\citeproctext]{ref-van2006making}
Van Oorschot, W., 2006. Making the difference in social europe: Deservingness perceptions among citizens of european welfare states. Journal of European social policy 16, 23--42.

\bibitem[\citeproctext]{ref-van2000should}
Van Oorschot, W., 2000. Who should get what, and why? On deservingness criteria and the conditionality of solidarity among the public. Policy \& Politics 28, 33--48.

\end{CSLReferences}

\section{Appendix}\label{appendix}

\subsection{Regression model for the fairness rankings}\label{regression-model-for-the-fairness-rankings}

We used a rank-ordered logit model (e.g., \citeproc{ref-hausman1987specifying}{Hausman and Ruud, 1987}) to estimate the effect of a change in the expected number of asylum seekers under a relocation scheme on the probability of ranking the scheme higher in the fairness ranking. This model allows probabilities to depend on the differences among all alternatives in a latent ``fairness'' variable. We assume the latent variable to be a function of the expected asylum seekers (scaled), \(AS\), and we allow the coefficients to vary by treatment. This variable can be formulated as:
\[
    \mu_{ij}^{\tau} = \alpha_{j}^{\tau} + \beta_{j}^{\tau} \text{AS},
\]
for all \(i\) and \(j\) alternatives, with \(\alpha_{j}^{\tau}\) and \(\beta_{j}^{\tau}\) set to zero for an arbitrary alternative and an arbitrary treatment condition \(\tau\). We selected the no-relocation scheme as the reference alternative, and the reference treatment group is the absolute-information group, so the model yields estimates of differences between control vs absolute and absolute vs relative information.

\subsection{Questionnaire}\label{questionnaire}

Welcome to this survey!

Thank you for your interest in this survey. This survey has been funded by the Joint Research Centre (JRC) of the European Commission and is implemented by Ipsos.

We are interested in your views about migration and asylum in the European Union. Please read all the information and answer the questions carefully. Different people will give different answers. It is your opinion that matters to us, do not hesitate to express it.
Your answers throughout this survey will be kept confidential. The European Commission will not receive any information from which they would be able to directly identify you. The data will be held for no longer than 5 years. The European Commission acts as data controller for all further data processing activities conducted by the European Commission on the research data delivered by Ipsos.

Participation in the survey is voluntary and you may withdraw consent at any time.

Before agreeing, please also read this information sheet {[}add hyperlink to information sheet{]}.

If you would like to read the Privacy Notice you can access it at \url{https://survey.ipsos.be/privacynotice_JRC_Refugees.pdf}

By agreeing to take part in this survey, you confirm that:

You have read the information about the survey (please click here {[}insert hyperlink to information sheet{]} to read the information sheet).

You are taking part in this survey by your own free will.

Do you agree to participate given the above conditions?

\begin{itemize}
\tightlist
\item
  Yes -- I have read the information above and agree to take part in the survey.
\item
  No -- I do not agree to take part in this survey. {[}screen out{]}
\end{itemize}

Base: all participants

D1. YEAR/MONTH. What is your date of birth?

\begin{itemize}
\tightlist
\item
  YEAR
\item
  \_1910 1910
\item
  \ldots{}
\item
  \_2015 2015
\item
  MONTH
\item
  \_1 January
\item
  \_2 February
\item
  \_3 March
\item
  \_4 April
\item
  \_5 May
\item
  \_6 June
\item
  \_7 July
\item
  \_8 August
\item
  \_9 September
\item
  \_10 October
\item
  \_11 November
\item
  \_12 December
\end{itemize}

{[}Standard Screener: DO NOT MODIFY OR TRANSLATE{]}

QUOTAGERANGE {[}Hidden{]}. Hidden Question - QUOTAGERANGE `this is a dummy question that will hold age breaks' for the quotas that should be defined by the PM; it CAN be edited and lines can be added to meet survey objectives.

\begin{itemize}
\tightlist
\item
  \_18\_24 `18-24',
\item
  \_25\_34 `25-34',
\item
  \_35\_44 `35-44'
\item
  \_45\_54 `45-54'
\item
  \_55\_64 `55-64'
\item
  \_65\_99 `65 and older'
\end{itemize}

{[}TERMINATE IF LESS THAN 18{]}

D2. {[}S{]} GENDER\_NONBINARY\_. Are you\ldots?

\begin{itemize}
\tightlist
\item
  \_1 Male
\item
  \_2 Female
\item
  \_3 Other
\item
  \_99 Prefer not to answer
\end{itemize}

D3. {[}S{]} In which region do you live?

1 Bulgaria See excel Region
2 France IIS standard screener - FRREGION5
3 Germany IIS standard screener - GERREGION1
4 Greece IIS standard screener - GRREGION1
5 Italy IIS standard screener - ITREGION1
6 Poland IIS standard screener - PLREGION1
7 Spain IIS standard screener - ESREGION2
8 Sweden IIS standard screener - SEREGION2
1. Prefer not to answer {[}hidden{]}\\
SCRIPTER: if D3=99: screenout

D4. {[}S{]} What is the highest level of education you have successfully completed (usually by obtaining a certificate or diploma)?

SCRIPTER: insert country-specific list Excel `Education'
SCRIPTER: add recode into ISCED (1 `Low', 2 `Mid', 3 `High', 99 `Unknown) and ISCED2 (1 `Low+Mid', 2 `High', 99 `Unknown)
1. Prefer not to answer

Base: All participants
D5. {[}S{]} Would you say that you live in a..?

\begin{enumerate}
\def\labelenumi{\arabic{enumi}.}
\tightlist
\item
  Rural area or village\\
\item
  Small or medium size town\\
\item
  Large town or city\\
\item
  Don't know (hidden)
\end{enumerate}

Q1. {[}Progressive grid{]} Would you say that most people can be trusted, or that you can't be too careful in dealing with people?\\
Please choose on a scale from 0 (You can't be too careful) to 10 (Most people can be trusted).

0
You can't be too careful 1 2 3 4 5 6 7 8 9 10
Most people can be trusted

Q2. {[}Progressive grid{]} How much do you trust your national government?\\
Please choose on a scale from 0 (Not at all) to 10 (Completely).

0
Not at all 1 2 3 4 5 6 7 8 9 10
Completely

Q3. {[}Progressive grid{]} How much do you trust the institutions of the European Union (EU)?\\
Please choose on a scale from 0 (Not at all) to 10 (Completely).

0
Not at all 1 2 3 4 5 6 7 8 9 10
Completely

Info1. Before moving on to the next block of questions, please note the differences between two terms: `asylum seekers' and `refugees':

SCRIPTER: insert as grid (check if readable on mobile)
Legal status Access to rights and services
Asylum seekers Application being processed:\\
Asylum seekers are individuals who have applied for asylum (protection) in another country, but their application is still being processed. Limited rights and access to services:\\
Asylum seekers may have limited access to rights and services depending on the laws and policies of the country they are in.
Refugees Officially recognised:\\
Refugees are individuals whose claim for international protection has been officially recognized according to the laws of the country they are seeking refuge in. Entitled to specific rights and protection:\\
Refugees are entitled to certain rights and protections under international law, including the right to work, access to education, and healthcare.

Q4. {[}S{]} (reverse scale for Random\_group=2) How much do you know about the migration and asylum system in the EU?

\begin{enumerate}
\def\labelenumi{\arabic{enumi}.}
\tightlist
\item
  I have no knowledge\\
\item
  I have basic knowledge\\
\item
  I have a good understanding\\
\item
  I have great expertise
\end{enumerate}

Q5. {[}S{]} How many asylum applications were received by {[}COUNTRY{]} in 2023?\\
Provide your best guess.

\begin{enumerate}
\def\labelenumi{\arabic{enumi}.}
\item
  Less than X1
\item
  Between X1 and X2
\item
  Between X2 and X3
\item
  Between X3 and X4
\item
  More than X4
\item
  I don't know
\end{enumerate}

Q6. {[}S{]} What proportion of asylum applications in {[}COUNTRY{]} was approved last year? Provide your best guess.

\begin{enumerate}
\def\labelenumi{\arabic{enumi}.}
\tightlist
\item
  0-10\%
\item
  11-20\%
\item
  21-30\%
\item
  31-40\%
\item
  41-50\%
\item
  51-60\%
\item
  61-70\%
\item
  71-80\%
\item
  81-90\%
\item
  91-100\%
\item
  I don't know
\end{enumerate}

Q7. {[}S{]} (reverse scale for Random\_group=2) Do you think that {[}COUNTRY{]} received more or fewer asylum requests than the EU average, last year?

\begin{enumerate}
\def\labelenumi{\arabic{enumi}.}
\tightlist
\item
  More asylum requests in {[}COUNTRY{]} than the EU average
\item
  Roughly the same
\item
  Fewer asylum requests in {[}COUNTRY{]} than the EU average
\item
  I don't know (fixed)
\end{enumerate}

Q8. {[}S{]} (reverse scale for Random\_group=2) Do you think that {[}COUNTRY{]} is faster or slower in processing asylum requests than the EU average?

\begin{enumerate}
\def\labelenumi{\arabic{enumi}.}
\tightlist
\item
  Faster in processing asylum requests than EU average
\item
  Roughly the same
\item
  Slower in processing asylum requests than EU average
\item
  I don't know (fixed)
\end{enumerate}

Q9. {[}S{]} (reverse scale for Random\_group=2) Are refugees in {[}COUNTRY{]} entitled to more rights or get access to more services than refugees in most other countries in the EU?

\begin{enumerate}
\def\labelenumi{\arabic{enumi}.}
\tightlist
\item
  More rights/services for refugees in {[}COUNTRY{]} than in most EU countries
\item
  Roughly the same
\item
  Fewer rights/services for refugees in {[}COUNTRY{]} than in most EU countries
\item
  I don't know (fixed)
\item
  Prefer not to answer (fixed)
\end{enumerate}

Q10. {[}Progressive grid{]} (reverse scale for Random\_group=2) In your opinion, how are asylum seekers currently distributed among the EU countries? Please answer on a scale from 0 to 10.

0
Unevenly distributed across the EU 1 2 3 4 5 6 7 8 9 10
Evenly distributed across the EU

Base: Random\_group=1
Q10v1. {[}Progressive grid{]} In your opinion, how are asylum seekers currently distributed among the EU countries? Please answer on a scale from 0 to 10.

0
Unevenly distributed across the EU 1 2 3 4 5 6 7 8 9 10
Evenly distributed across the EU

Base: Random\_group=2\\
Q10v2. {[}Progressive grid{]} In your opinion, how are asylum seekers currently distributed among the EU countries? Please answer on a scale from 0 to 10.

0
Evenly distributed across the EU 1 2 3 4 5 6 7 8 9 10
Unevenly distributed across the EU

Q11. {[}Progressive grid{]} How important is it that the European Union system for hosting asylum seekers and refugees ensures that \ldots?

Rows (randomize):
1. \ldots the cost related to hosting asylum seekers or refugees is shared among EU countries
2. \ldots the number of asylum seekers or refugees hosted is not too different among EU countries
3. \ldots asylum seekers or refugees have access to the same rights across all EU countries
4. \ldots asylum seekers or refugees are treated in the same way in all EU countries
5. \ldots asylum applications are processed quickly

Columns:

\begin{enumerate}
\def\labelenumi{\arabic{enumi}.}
\tightlist
\item
  0 - Not at all important\\
\item
  1
\item
  2
\item
  3
\item
  4
\item
  5
\item
  6
\item
  7
\item
  8
\item
  9
  10 - Very important
\end{enumerate}

Q12. Of the five considerations listed in the previous screen, which is the most important to you?

\begin{enumerate}
\def\labelenumi{\arabic{enumi}.}
\tightlist
\item
  the cost related to hosting asylum seekers or refugees is shared among EU countries\\
\item
  the number of asylum seekers or refugees hosted is not too different among EU countries\\
\item
  asylum seekers or refugees have access to the same rights across all EU countries\\
\item
  asylum seekers or refugees are treated in the same way in all EU countries\\
\item
  asylum applications are processed quickly
\end{enumerate}

Q13. {[}S{]} (reverse scale for Random\_group=2) Do you believe that {[}COUNTRY{]} is currently hosting a fair share of the total asylum seekers in the European Union?

\begin{enumerate}
\def\labelenumi{\arabic{enumi}.}
\tightlist
\item
  Yes
\item
  No, the share is too high\\
\item
  No, the share is too low
\item
  I don't know (fixed)
\item
  Prefer not to answer (fixed)
  PROG: create 3 equal groups: create hidden variable `relocation\_treatment', according to least fill quota per country. Assign each participant to one of these 3 groups. For group 2 and 3 extra info is to be shown in Q13:
\item
  Control
\item
  Absolute
\item
  Relative
\end{enumerate}

Q14. {[}S{]} Below you see three scenarios describing how asylum applicants could be distributed across the EU. \textless IF TREATED: The numbers are based on national statistics.\textgreater{} In your view, what would be a fair way to establish the number of asylum applicants per country? Please rank the three options from the fairest to the least fair one.\\

\begin{enumerate}
\def\labelenumi{\arabic{enumi}.}
\tightlist
\item
  No relocation: Applications should be handled in the country of first entry (that is, asylum seekers are required to submit their asylum application in the European country in which they initially arrive).\\
  PROG: show for relocation\_treatment =2: `This would mean approximately \textless Bold: X1\textgreater{} asylum applications allocated to {[}COUNTRY{]} on average per year.'\\
  PROG: show for relocation\_treatment =3: `This would mean approximately \textless Bold: Y1\textgreater{} applications allocated to {[}COUNTRY{]} per 100,000 of its citizens on average per year.'
\item
  Proportional to population: The number of applications should be proportional to the country's population.\\
  PROG: show for relocation\_treatment =2: `This would mean approximately \textless Bold: X2\textgreater asylum applications allocated to {[}COUNTRY{]} on average per year.'\\
  PROG: show for relocation\_treatment =3: `This would mean approximately \textless Bold: Y2\textgreater{} applications allocated to {[}COUNTRY{]} per 100,000 of its citizens on average per year.'
\item
  Proportional to the country's economy: The number of applications should be proportional to the country's GDP.\\
  PROG: show for relocation\_treatment =2: `This would mean approximately \textless Bold: X3\textgreater asylum applications allocated to {[}COUNTRY{]} on average per year.'\\
  PROG: show for relocation\_treatment =3: `This would mean approximately \textless Bold: Y3\textgreater{} applications allocated to {[}COUNTRY{]} per 100,000 of its citizens on average per year.'
  PROG: replace value and by country specific numbers (see xls file)
  Scripter: Please keep Q14 and Q15 on same page.
\end{enumerate}

Q15. Could you please briefly explain why you ranked the options the way you did? What influenced your choice?
PROG: create 3 random groups (completely random, no least fill quota): create hidden variable `scenario'. Assign each participant to one of these 3 groups. Different info shown in Q16 depending on group:
1. fair
2. unfair
3. unfair\_more

Q16. {[}SLIDER{]} Imagine a hypothetical situation where \textless if Scenario=1 or 2, show: `200,000' if Scenario=3, show `300,000'\textgreater{} refugees from a non-EU country need to be resettled across three EU countries \textless PROG: show: `({[}countryX{]}, {[}countryY{]}, and {[}COUNTRY{]})\textgreater. The refugees are initially distributed as follows:
PROG: show for Treat Scenario=1 (Fair Distribution):\\
* {[}countryX{]}: 55,000 refugees
* {[}countryY{]}: 45,000 refugees

PROG: show for Treat Scenario=2 (Less Fair Distribution):\\
* {[}countryX{]}: 85,000 refugees
* {[}countryY{]}: 15,000 refugees

PROG: show for Treat Scenario=3 (Less Fair, More Refugees)
* {[}countryX{]}: 135,000 refugees
* {[}countryY{]}: 65,000 refugees

Based on the scenario: how many of the remaining 100,000 refugees would you be willing to accept in {[}COUNTRY{]}? If you select a number less than 100,000, the remaining amount will be evenly distributed between the other two countries.

{[}INSERT SLIDER ranging from 0 to 100,000, in steps of 5000, show value of selected point{]}

PROG: replace {[}countryX{]} and {[}countryY{]} with random pairs from \{Romania, Belgium, Austria, Czech Republic\}, store the country shown in variable countryX and countryY:

Pair CountryX CountryY
1 Romania Belgium
2 Romania Austria
3 Romania Czech Republic
4 Belgium Romania
5 Belgium Austria
6 Belgium Czech Republic
7 Austria Romania
8 Austria Belgium
9 Austria Czech Republic
10 Czech Republic Romania
11 Czech Republic Belgium
12 Czech Republic Austria

Base: if pilot=Yes
Q16b. {[}O{]} When responding to the previous question, was there anything unclear or difficult to understand? If so, what? Please also share any suggestions you may have to improve the question.\,
\_\_\_

\begin{enumerate}
\def\labelenumi{\arabic{enumi}.}
\setcounter{enumi}{97}
\tightlist
\item
  Everything was clear
\end{enumerate}

Q17. {[}S{]} (reverse scale for Random\_group=2) In your opinion, how important is it that EU countries hosting more asylum seekers receive financial support from other EU countries?

\begin{enumerate}
\def\labelenumi{\arabic{enumi}.}
\tightlist
\item
  Very important
\item
  Somewhat important
\item
  Not important at all
\item
  I don't know (fixed)
\end{enumerate}

Q18. {[}S{]} (randomize) To the best of your knowledge, who handles asylum applications in the EU?

\begin{enumerate}
\def\labelenumi{\arabic{enumi}.}
\tightlist
\item
  EU authorities
\item
  National authorities\\
\item
  I don't know (fixed)
\end{enumerate}

Q19. {[}Progressive grid{]} People have different opinions about the reasons to seek asylum. In your opinion, to what extent do the following reasons justify that a person seeks asylum in {[}COUNTRY{]}?

Rows (randomize, keep 3+4 together):
1. Escaping persecution due to nationality, ethnic origin, political opinion, or sexual orientation \\
2. Escaping violence, conflict or war \\
3. Seeking employment \\
4. Seeking social welfare \\
Columns (reverse scale for Random\_group=2):\\
1. Completely justified
2. Justified
3. Not at all justified
98. I don't know (fixed)
99. Prefer not to answer (fixed)

Q20. {[}S{]} (randomize 1-4, same order as Q19) In your opinion, what is the most common reason for a person to seek asylum in {[}COUNTRY{]}?

\begin{enumerate}
\def\labelenumi{\arabic{enumi}.}
\tightlist
\item
  Escaping persecution due to nationality, ethnic origin, political opinion, or sexual orientation
\item
  Escaping violence, conflict or war \\
\item
  Seeking employment
\item
  Seeking social welfare
\item
  Other, please specify {[}O{]} (fixed)
\item
  I don't know (fixed)
\item
  Prefer not to answer (fixed)
\end{enumerate}

Q\_Check. {[}S{]} Please read the following statements. Please be so kind and let us know that you're paying attention and reading all the questions carefully: please choose the second statement below, regardless if you agree with it or not. The most desirable characteristic of an asylum system is that\ldots{}

\begin{enumerate}
\def\labelenumi{\arabic{enumi}.}
\tightlist
\item
  it is cost-effective\\
\item
  decisions are free from prejudice
\item
  decisions are made quickly
\item
  the interests of EU Member States are safeguarded
\item
  the rights of the applicants are observed
\end{enumerate}

PROG: if Q\_Check \textless\textgreater{} 2 à Screen out

Base: all participants (filled during data processing, one of the 2 below versions is shown to participants)
PROG: Please combine Q21v1 and Q21v2 in 1 variable Q21, where the labels of Q21 are used (so copy Q21v1 as is and recode Q21v2 to reverse the codes 0=10, 1=9, \ldots, 10=0).
Q21. {[}Progressive grid{]} (reverse scale for Random\_group=2) Do you think that asylum decisions taken by immigration authorities in {[}COUNTRY{]} are usually based on evidence and legal standards or are usually influenced by political, cultural, or personal biases?\\
Rate on a scale of 0 to 10, where 0 indicates all decisions are made solely based on verifiable evidence and legal standards, and where 10 suggests decisions are heavily influenced by various biases.

0
Asylum decisions are made solely based on verifiable evidence and legal standards 1 2 3 4 5 6 7 8 9 10
Asylum decisions are heavily influenced by political, cultural, or personal biases

Base: Random\_group=1
Q21v1. {[}Progressive grid{]} Do you think that asylum decisions taken by immigration authorities in {[}COUNTRY{]} are usually based on evidence and legal standards or are usually influenced by political, cultural, or personal biases?\\
Rate on a scale of 0 to 10, where 0 indicates all decisions are made solely based on verifiable evidence and legal standards, and where 10 suggests decisions are heavily influenced by various biases.

0
Asylum decisions are made solely based on verifiable evidence and legal standards 1 2 3 4 5 6 7 8 9 10
Asylum decisions are heavily influenced by political, cultural, or personal biases

Base: Random\_group=2
Q21v2. {[}Progressive grid{]} Do you think that asylum decisions taken by immigration authorities in {[}COUNTRY{]} are usually based on evidence and legal standards or are usually influenced by political, cultural, or personal biases?\\
Rate on a scale of 0 to 10, where 0 indicates all decisions are heavily influenced by various biases, and where 10 suggests decisions are made solely based on verifiable evidence and legal standards.

0
Asylum decisions are heavily influenced by political, cultural, or personal biases 1 2 3 4 5 6 7 8 9 10
Asylum decisions are made solely based on verifiable evidence and legal standards

Base: all participants (filled during data processing, one of the 2 below versions is shown to participants)\\
PROG: Please combine Q22v1 and Q22v2 in 1 variable Q22, where the labels of Q22 are used (so copy Q22v1 as is and recode Q22v2 to reverse the codes 0=10, 1=9, \ldots, 10=0).
Q22. {[}S{]} (reverse scale for Random\_group=2) Suppose that some asylum applications submitted in {[}COUNTRY{]} contain false or misleading information. How confident are you in the ability of your country's asylum system to identify these asylum applications?\\
Rate it on a scale of 0 to 10, where 0 indicates Not confident at all (the system fails to recognize a significant number of false applications) and 10 indicates that Completely confident (the system detects false applications effectively).

0
Not confident at all 1 2 3 4 5 6 7 8 9 10
Completely confident

Base: Random\_group=1
Q22v1. {[}S{]} Suppose that some asylum applications submitted in {[}COUNTRY{]} contain false or misleading information. How confident are you in the ability of your country's asylum system to identify these asylum applications?\\
Rate it on a scale of 0 to 10, where 0 indicates Not confident at all (the system fails to recognize a significant number of false applications) and 10 indicates Completely confident (the system detects false applications effectively).

0
Not confident at all 1 2 3 4 5 6 7 8 9 10
Completely confident

Base: Random\_group=2
Q22v2. {[}S{]} Suppose that some asylum applications submitted in {[}COUNTRY{]} contain false or misleading information. How confident are you in the ability of your country's asylum system to identify these asylum applications?\\
Rate it on a scale of 0 to 10, where 0 indicates Completely confident (the system detects false applications effectively) and 10 indicates Not confident at all (the system fails to recognize a significant number of false applications).

0
Completely confident 1 2 3 4 5 6 7 8 9 10
Not confident at all

Q23. {[}S{]} In your opinion, should asylum seekers have the right to decide on their final country of destination in the EU?

\begin{enumerate}
\def\labelenumi{\arabic{enumi}.}
\tightlist
\item
  Yes
\item
  No
\item
  I don't know
\item
  Prefer not to answer
\end{enumerate}

Base: all participants
Q24. {[}O{]} Do these terms: migrant, asylum seeker, refugee, mean the same?

If not, please briefly describe how you understand the differences between them.

Info2. Please pay attention to the following definition of migrants before continuing the survey:

Migrants: We refer to migrants as people who are citizens of a country outside the EU and who move to one EU country. They might move for various reasons: escaping danger, seeking work or a better life. Asylum seekers and refugees are two categories of migrants, but not all migrants are asylum seekers or refugees.
SCRIPTER: insert country specific image (translations will be included):

Q25. {[}Q{]} What proportion of migrants from outside the EU that are currently in {[}COUNTRY{]} are asylum seekers?
Provide your best guess.

\begin{enumerate}
\def\labelenumi{\arabic{enumi}.}
\tightlist
\item
  0-10\%
\item
  11-20\%
\item
  21-30\%
\item
  31-40\%
\item
  41-50\%
\item
  51-60\%
\item
  61-70\%
\item
  71-80\%
\item
  81-90\%
\item
  91-100\%
\item
  I don't know
\end{enumerate}

Q26. {[}Q{]} What proportion of the residents of {[}COUNTRY{]} are migrants from outside the EU?
Provide your best guess.

\begin{enumerate}
\def\labelenumi{\arabic{enumi}.}
\tightlist
\item
  0-10\%
\item
  11-20\%
\item
  21-30\%
\item
  31-40\%
\item
  41-50\%
\item
  51-60\%
\item
  61-70\%
\item
  71-80\%
\item
  81-90\%
\item
  91-100\%
\item
  I don't know
\end{enumerate}

PROG: create 2 random groups (completely random, no least fill quota): create hidden variable `Q27\_place'. Assign each participant to one of these 2 groups. Q27 to be shown here or after Q33. Please combine Q27a and Q27b in 1 variable Q27.
1. After Q26
2. After Q33

Base: IF Q27\_place=1
Q27a. {[}SLIDER{]} Among migrants from outside the EU staying in {[}COUNTRY{]}, how many do you think are staying legally?

{[}INSERT SLIDER with left and right label, ranging from 0\% to 100\%, show value of selected point{]}
0\% staying legally 100\% staying legally

\begin{enumerate}
\def\labelenumi{\arabic{enumi}.}
\setcounter{enumi}{97}
\tightlist
\item
  I don't know
\item
  Prefer not to answer
\end{enumerate}

Q28. {[}Progressive grid{]} Do you support or oppose that migrants from outside the EU have access to the following services in {[}COUNTRY{]}?

Rows (randomize):
1. Public healthcare
2. Public education for minors
3. Social assistance
4. Public housing
5. Access to the labour market

Columns (reverse scale for Random\_group=2):
1. Strongly oppose\\
2. Mildly oppose\\
3. Neutral\\
4. Mildly support
5. Strongly support
98. I don't know (fixed)
99. Prefer not to answer (fixed)

Q29. {[}Progressive grid{]} Do you support or oppose that migrants from outside the EU receive the following financial benefits in {[}COUNTRY{]}?

Rows (randomize):
1. Family allowance
2. Unemployment benefits
3. Disability allowance

Columns (reverse scale for Random\_group=2):\\
1. Strongly oppose\\
2. Mildly oppose\\
3. Neutral\\
4. Mildly support
5. Strongly support
98. I don't know (fixed)
99. Prefer not to answer (fixed)

Q30. {[}S{]} (reverse scale for Random\_group=2) In your opinion, how does the access to public services (healthcare, education, housing etc.) of migrants from outside the EU compare to that of citizens in {[}COUNTRY{]}?

\begin{enumerate}
\def\labelenumi{\arabic{enumi}.}
\tightlist
\item
  It is much better for migrants
\item
  It is somewhat better for migrants
\item
  It is about the same
\item
  It is somewhat worse for migrants
\item
  It is much worse for migrants
\item
  I don't know (fixed)
\item
  Prefer not to answer (fixed)
\end{enumerate}

Q31. {[}S{]} (reverse scale for Random\_group=2) In your opinion, how does the level of financial benefits provided to migrants from outside the EU compare to that of citizens in {[}COUNTRY{]}?

\begin{enumerate}
\def\labelenumi{\arabic{enumi}.}
\tightlist
\item
  It is much better for migrants\\
\item
  It is somewhat better for migrants\\
\item
  It is about the same
\item
  It is somewhat worse for migrants\\
\item
  It is much worse for migrants
\item
  I don't know (fixed)
\item
  Prefer not to answer (fixed)
\end{enumerate}

SCRIPTER: Randomise order in which Q32 and Q33 are shown. Capture order in variable

Q32. {[}S{]} (reverse scale for Random\_group=2) To what extent are you concerned that fewer benefits and services will be available to citizens, if benefits and services are provided to migrants from outside the EU?

\begin{enumerate}
\def\labelenumi{\arabic{enumi}.}
\tightlist
\item
  Very concerned
\item
  Somewhat concerned
\item
  Not really concerned
\item
  Not at all concerned
\item
  I don't know (fixed)
\item
  Prefer not to answer (fixed)
\end{enumerate}

Q33. {[}S{]} (reverse scale for Random\_group=2) In your opinion, how much do migrants from outside the EU contribute to the economy in {[}COUNTRY{]} (for instance, by paying taxes)?

\begin{enumerate}
\def\labelenumi{\arabic{enumi}.}
\tightlist
\item
  More than citizens
\item
  As much as citizens\\
\item
  Less than citizens\\
\item
  I don't know (fixed)
\item
  Prefer not to answer (fixed)
\end{enumerate}

Base: IF Q27\_place=2
Q27b. {[}SLIDER{]} Among migrants from outside the EU staying in {[}COUNTRY{]}, how many do you think are staying legally?

{[}INSERT SLIDER with left and right label, ranging from 0\% to 100\%, show value of selected point{]}
0\% staying legally 100\% staying legally

\begin{enumerate}
\def\labelenumi{\arabic{enumi}.}
\setcounter{enumi}{97}
\tightlist
\item
  I don't know
\item
  Prefer not to answer
\end{enumerate}

Q34. {[}M{]} (randomize) Where do you get most of your news on migration matters?\\
Select all that apply.

\begin{enumerate}
\def\labelenumi{\arabic{enumi}.}
\tightlist
\item
  Television
\item
  Websites
\item
  Online Social Networks
\item
  Radio
\item
  Printed press
\item
  Other (fixed)
\item
  I do not look for news on migration matters (fixed) (exclusive) \\
\item
  I don't know (fixed) (exclusive)
\end{enumerate}

Q35. {[}S{]} (reverse scale for Random\_group=2) How often do you hear news about migration?

\begin{enumerate}
\def\labelenumi{\arabic{enumi}.}
\tightlist
\item
  Very often\\
\item
  Not too often\\
\item
  Not often\\
\item
  Never
\end{enumerate}

Q36. {[}S{]} (reverse scale for Random\_group=2) In your opinion, are migrants portrayed accurately by the media in {[}COUNTRY{]}?

\begin{enumerate}
\def\labelenumi{\arabic{enumi}.}
\tightlist
\item
  No, they are portrayed too positively
\item
  Yes, they are portrayed accurately
\item
  No, they are portrayed too negatively
\item
  I don't know (fixed)
\end{enumerate}

Info3. The next few questions may be considered more sensitive in nature. This includes questions on ethnicity, political opinions, and income. We are asking them because they will help us to better understand your opinion.

We want to remind you that for each question a `Prefer not to answer' option is available for you to select. Participation is always voluntary, and your responses are used for research purposes only.\\
The European Commission will not receive any information from which they would be able to directly identify you. The data will be held for no longer than 5 years. The European Commission acts as data controller for all further data processing activities conducted by the European Commission on the research data delivered by Ipsos.

Q38. {[}S{]} Which of these descriptions best describes your current situation?

\begin{enumerate}
\def\labelenumi{\arabic{enumi}.}
\tightlist
\item
  In paid work
\item
  In education
\item
  Unemployed and actively looking for a job
\item
  Unemployed, wanting a job, but not actively looking for a job \textless{} LMS\_inactive\textgreater{}
\item
  Permanently sick or disabled
\item
  Retired
\item
  In community or military service
\item
  Doing housework, looking after children or other persons
\item
  Other
\item
  Prefer not to answer (exclusive)
\end{enumerate}

Q39. {[}S{]} In what country were you born?

{[}INSERT DROPDOWN list country{]}
999. Prefer not to answer

Q40. {[}S{]} How many years have you been living in {[}COUNTRY{]}?

\_\_\_\_\_\_ {[}NUMERIC, valid range: 0-participant age{]}
99. Prefer not to answer

Q41. {[}S{]} In what country was your father born?

{[}INSERT DROPDOWN list country{]}
998. I don't know
999. Prefer not to answer

Q42. {[}S{]} In what country was your mother born?

{[}INSERT DROPDOWN list country{]}
998. I don't know
999. Prefer not to answer

Q43. {[}S{]} How many children do you have?

\begin{enumerate}
\def\labelenumi{\arabic{enumi}.}
\tightlist
\item
  I do not have children:\\
\item
  1
\item
  2\\
\item
  3\\
\item
  4\\
\item
  5 or more\\
\item
  Prefer not to answer
\end{enumerate}

Q44. {[}S{]} What is your relationship status?

\begin{enumerate}
\def\labelenumi{\arabic{enumi}.}
\tightlist
\item
  Single
\item
  Married
\item
  Legally separated or divorced
\item
  Widowed
\item
  Prefer not to answer
\end{enumerate}

Q45. {[}S{]} How many people can you count on supporting you in case of need?

\begin{enumerate}
\def\labelenumi{\arabic{enumi}.}
\tightlist
\item
  None
\item
  One\\
\item
  Two
\item
  Three
\item
  Four to six
\item
  Seven to nine
\item
  Ten or more
\item
  Prefer not to answer
\end{enumerate}

Q46. {[}prorgessive grid{]} We would like to understand what you value.\\
Please read each description and respond how much the description is or is not like you.\,

Rows:
1. I believe I should always show respect to my parents and to older people. It is important to me to be obedient. \\
2. Religious belief is important to me. I try hard to do what my religion requires. \\
3. It's very important to me to help the people around me. I want to care for their well-being. \\
4. I think it is important that every person in the world is treated equally. I believe everyone should have equal opportunities in life.
5. I think it's important to be interested in things. I like to be curious and to try to understand all sorts of things. \\
6. I like to take risks. I am always looking for adventures. \\
7. I seek every chance to have fun. It is important to me to do things that give me pleasure. \\
8. Being very successful is important to me. I like to impress other people. \\
9. It is important to me to be in charge and tell others what to do. I want people to do what I say. \\
10. It is important to me that things are organised and clean. I really do not like things to be a mess.

Columns\\
1. Very much like me
2. Like me
3. Somewhat like me
4. A little like me
5. Not like me
6. Not at all like me
99. Prefer not to answer

Q47. {[}S{]} PROG: If country=1,4, 6, 8 Show: Could you please indicate your household's monthly income (that is, after income taxes have been paid)?
PROG: If country=2, 3, 5, 7 Show: Could you please indicate your household's annual income (that is, after income taxes have been paid)?

Your total household income includes your own income plus the incomes of all household members who live together with you. The total income includes incomes from jobs, pensions, social security, interest, dividends, capital gains claimed, profits from businesses, unemployment payments, and all other money your received.

Scripter: insert values from country-specific list Excel `Q47\_Income'

\begin{enumerate}
\def\labelenumi{\arabic{enumi}.}
\tightlist
\item
  Less than {[}insert value1{]} {[}insert currency{]}
\item
  Between {[}insert value1{]} {[}insert currency{]} and {[}insert value2 minus 1{]} {[}insert currency{]}
\item
  Between {[}insert value2{]} {[}insert currency{]} and {[}insert value3 minus 1{]} {[}insert currency{]}
\item
  Between {[}insert value3{]} {[}insert currency{]} and {[}insert value4 minus 1{]} {[}insert currency{]}
\item
  Between {[}insert value4{]} {[}insert currency{]} and {[}insert value5 minus 1{]} {[}insert currency{]}
\item
  Between {[}insert value5{]} {[}insert currency{]} and {[}insert value6 minus 1{]} {[}insert currency{]}
\item
  Between {[}insert value6{]} {[}insert currency{]} and {[}insert value7 minus 1{]} {[}insert currency{]}
\item
  Between {[}insert value7{]} {[}insert currency{]} and {[}insert value8 minus 1{]} {[}insert currency{]}
\item
  Between {[}insert value8{]} {[}insert currency{]} and {[}insert value9{]} {[}insert currency{]}
\item
  Higher than {[}insert value9{]} {[}insert currency{]}
\item
  I don't know
\item
  Prefer not to answer
  Scripter: recode Q47 into Q47\_quota:\\
\item
  `Low` IF Q47 = 1, 2, 3, 4
\item
  `Mid' IF Q47 = 5, 6, 7
\item
  `High' IF Q47 = 8, 9, 10
\item
  `Missing' IF Q47 = 98 or 99
\end{enumerate}

Q48. {[}S{]} How does your income compare to that of the people around you?

\begin{enumerate}
\def\labelenumi{\arabic{enumi}.}
\tightlist
\item
  I earn a somewhat lower income
\item
  I earn an income similar to that of the people around me
\item
  I earn a somewhat higher income
\item
  I don't know
\item
  Prefer not to answer
\end{enumerate}

Q49. {[}S{]} In political matters people talk of `the left' and `the right'. How would you place your views on this scale?

0
Left 1 2 3 4 5 6 7 8 9 10
Right

\begin{enumerate}
\def\labelenumi{\arabic{enumi}.}
\setcounter{enumi}{98}
\tightlist
\item
  Prefer not to answer
\end{enumerate}

Q37. {[}S{]} (reverse scale for Random\_group=2) Overall, how satisfied are you with the life you lead?

\begin{enumerate}
\def\labelenumi{\arabic{enumi}.}
\tightlist
\item
  Very satisfied
\item
  Fairly satisfied
\item
  Not very satisfied
\item
  Not at all satisfied
\item
  Prefer not to answer
\end{enumerate}

Q50. For your information, you were asked to indicate whether you think that {[}COUNTRY{]} received more or fewer asylum requests than the EU average last year. The EU average for last year was 38,867. With {[}insert number from excel Q7\_asylum requests column B{]} asylum requests last year, {[}COUNTRY{]} received {[}insert fewer/more from excel Q7\_asylum requests column D{]} asylum requests than the EU average.

\subsection{Additional Tables}\label{additional-tables}

\begin{table}
\centering
\caption{\label{tab:assignment-summary}Random assignment of respondents, by country}
\centering
\begin{tabu} to \linewidth {>{}l>{\raggedleft}X>{\raggedleft}X>{\raggedleft}X>{\raggedleft}X}
\toprule
  & Control & Absolute & Relative & Total\\
\midrule
\textbf{Bulgaria} & 332 & 337 & 332 & 1 001\\
\textbf{France} & 339 & 339 & 337 & 1 015\\
\textbf{Germany} & 343 & 337 & 344 & 1 024\\
\textbf{Greece} & 337 & 332 & 331 & 1 000\\
\textbf{Italy} & 337 & 343 & 345 & 1 025\\
\textbf{Poland} & 340 & 335 & 353 & 1 028\\
\textbf{Spain} & 343 & 350 & 334 & 1 027\\
\textbf{Sweden} & 332 & 331 & 339 & 1 002\\
\addlinespace
\textbf{Total} & 2 703 & 2 704 & 2 715 & 8 122\\
\bottomrule
\end{tabu}
\end{table}

\begin{figure}

{\centering \includegraphics[width=0.75\textwidth]{/Users/mrb/Research/migrant_relocation_fairness/draft/jebo-rev1/main_files/figure-latex/rologit-base-1} 

}

\caption[Rank-ordered logit baseline specification.]{Estimates and 95\% CIs from a rank-ordered logit regression on the fairness rankings. $\beta_{AS}$ represents the effect of a SD change in expected asylum seekers on the probability of ranking an allocation scheme as fairer. $\alpha_{j}$ represents the effect of a proportional allocation mechanism $j$ relative to the no relocation alternative (omitted). This shows that only a large increase in expected asylum seekers can fully offset the baseline perception of fairness for proportional systems.}\label{fig:rologit-base}
\end{figure}

\end{document}
